\ParallelTitle{Keywords and concepts}{关键词与概念}
\ParallelText{lookup-intro}
 {Lookup terms below lead to one canonical concept. Page pointers are generated from the notes, not typed manually.}
 {下面的检索词指向同一个规范概念条目。页码由笔记自动生成，并非手工输入。}
\ParallelText{alias-mean-speed}
 {\textbf{Mean speed.} See \ParallelReference{concept-speed}{Average speed}.}
 {\textbf{平均速率的别称。}见\ParallelReference{concept-speed}{平均速率}。}
\ParallelText{keyword-units}
 {\textbf{Units.} See \ParallelReference{concept-speed}{Average speed} for a distance/time example.}
 {\textbf{单位。}距离与时间的例子见\ParallelReference{concept-speed}{平均速率}。}
\ParallelSection{concept-speed}{Average speed}{平均速率}
\ParallelText{recall-speed}
 {Divide total path length by positive elapsed time. With metres and seconds, the result is in metres per second.}
 {用总路程除以大于零的总时间。若路程用米、时间用秒，结果的单位就是米每秒。}
\ParallelText{check-speed}
 {\textbf{Check:} multiply the result by elapsed time to recover the distance. This checks arithmetic, not measurement accuracy.}
 {\textbf{核验：}将结果乘以总时间，反算路程。这能检查计算，但不能保证测量准确。}
\ParallelText{qr-notes-speed}
 {\hyperref[notes-sample-concept]{Notes, section~\ref*{notes-sample-concept}, p.~\pageref*{notes-sample-concept}}}
 {\hyperref[notes-sample-concept]{笔记第\ref*{notes-sample-concept}节，第\pageref*{notes-sample-concept}页}}
